\documentclass{article} %210 mm × 297 mm
\usepackage[utf8]{inputenc}
\usepackage[a4paper,left=1.5cm, right=1.5cm, top=2cm, bottom=2cm]{geometry}
\usepackage{graphicx}
%\usepackage{blindtext}
\usepackage{multirow}
\usepackage{mhchem}
\usepackage{url}
\usepackage{subfigure}
\usepackage{amsfonts,latexsym} % para tener disponibilidad de diversos simbolos
\usepackage{enumerate}
%\usepackage{booktabs}
\usepackage{float}
%\usepackage{threeparttable}
\usepackage{array,colortbl}
%\usepackage{ifpdf}
%\usepackage{rotating}
\usepackage{cite}
\usepackage{stfloats}
\usepackage{url}
\usepackage{listings}
\usepackage{fancyhdr}

\usepackage{color}
%\usepackage{hyperref}
%\usepackage{wrapfig}
\usepackage{array}
%\usepackage{multirow}
%\usepackage{adjustbox}
%\usepackage{nccmath}
\usepackage[dvipsnames]{xcolor}


\newcommand{\titulo}{Präsenzübungsblatt 5; MfN II.}
\newcommand{\nombre}{Liliana Sanfilippo}
\newcommand{\FCE}{\textbf{SoSe 2025}}

 

\usepackage{fancyhdr}
\pagestyle{fancy}
\fancyhead{}
\fancyfoot{}
\fancyhead[LO]{\small\nombre}
\fancyhead[RE]{\small\titulo}
\fancyhead[LO]{\small\FCE}
\fancyfoot[RO,LE] {\thepage}

\setcounter{page}{1} %Número de la página, la editorial lo cambiará



\usepackage[onehalfspacing]{setspace}
\begin{document}


\begin{tabular}{p{12cm}}
{{\Large\bf\titulo}}\\
\vspace{0.2cm}\\
 {\large\nombre}\\
\end{tabular}
\vspace{0.2cm}\\
\section*{Besprechung 3.1 D}
\[
D = \begin{pmatrix}
2-i&0&0\\
i&-3&1-i\\
1&0&2+i
\end{pmatrix}
\]
wenn $\lambda_2 = -3$
\[
D+3I = \begin{pmatrix}
2-i+3&0&0\\
i&0&1-i\\
1&0&5+i
\end{pmatrix} \quad und \quad
E_D(\lambda_2) = ker \begin{pmatrix}
2-i+3&0&0\\
i&0&1-i\\
1&0&5+i
\end{pmatrix}
\]
Kern hat Dimension 1 und deswegen kann man das direkt ablesen 
\[
\Rightarrow \begin{pmatrix}
0 \\ 1 \\ 0
\end{pmatrix}
\]
wenn $\lambda_1 = 2-i$
\[
E_D(\lambda_1) = ker\begin{pmatrix}
0&0&0\\
i&-5+i&1-i\\
1&0&2i
\end{pmatrix} 
\]
\begin{align*}
    0 & = ix_2 + x_2 (-5 + i) + x_3 (1-i) & (i)\\
    0 & = x_1 + x_3 \cdot 2i  \\
    x_1 & = -x_3 \cdot 2i & (ii) 
\end{align*}
(ii) in (i) eingesetzt 
\[
2 x_3 + x_2 (-5 +i) + x_3 - i x_3 = 0
\]
\[x_2 (-5 +i) = x_3 (-3 +i) \]
\[x_2 = x_3 (\frac{-3+i}{-5+i}) \Rightarrow EV zu \lambda_1 ist x_3 \cdot \begin{pmatrix} -2i \\ \frac{-3+i}{-5+i} \\ 1 
\end{pmatrix} = x_3 \cdot \begin{pmatrix} -2i \\ \frac{8+i}{13} \\ 1 
\end{pmatrix}
\]
Wenn $\lambda_3 = 2+i$ \\
Hier bieten sich Zeilenumformungen an 
\[
E_D(\lambda_3) = ker\begin{pmatrix}
-2i&0&0\\
i&-5-i&1-i\\
1&0&0
\end{pmatrix} = ker\begin{pmatrix}
0&0&0\\
i&-5-i&1-i\\
1&0&0
\end{pmatrix}
\]
Man sieht an der ersten Zeile (nur nullen), dass die erste Zahl im Vektor eine Null sein muss. \\
Aus der dritten Zeile folgt: für $\begin{pmatrix} x_1 \\ x_2 \\ x_3 \end{pmatrix} \in ker \Rightarrow x_1 = 0$ \\
\[\Rightarrow (-5-i) \cdot x_2 + x_3 (1-i) = 0\]
\[\Rightarrow \begin{pmatrix} 0 \\ -(1-i) \\ -5-1 \end{pmatrix} \in ker \Rightarrow <\begin{pmatrix} 0 \\ -(1-i) \\ -5-1 \end{pmatrix} >  sind die EV zu \lambda_3 \]

Wenn die algebraische Vielfachheit 1 ist kann die geometrische auch nur 1 sein, weil die g immer kleiner gleich die a ist. 
\section*{Besprechung 3.1 E}
\[
 E = \begin{pmatrix}
 1 &-2&1 \\
 0&1&3\\
 0&-1&-2
 \end{pmatrix}
\]
\[
 E-\lambda I = \begin{pmatrix}
 1 -\lambda  &-2&1 \\
 0&1-\lambda &3\\
 0&-1&-2-\lambda 
 \end{pmatrix}
\]
\[
\chi_{\lambda} = det(E-\lambda I) 
\]
Laplace nach der ersten Zeile. 
\begin{align*}
    \chi_{\lambda} & = det(E-\lambda I) \\
    & = (1-\lambda) \cdot det \begin{pmatrix} 1- \lambda & 3 \\ -1 & -2- \lambda \end{pmatrix} \\
     & = (1-\lambda) \cdot (1+t+t^2) 
\end{align*}
\[
\lambda_1 = 1, \lambda_2 = -\frac{1}{2} + \frac{1}{2} \cdot \sqrt{3}i, \lambda_3 = -\frac{1}{2} - \frac{1}{2} \cdot \sqrt{3}i
\]
\[
E v_E(1) = ker \begin{pmatrix}
0 &-2&1 \\
 0&0&3\\
 0&-1&-3 
 \end{pmatrix} = < \begin{pmatrix} 1 \\ 0 \\ 0 \end{pmatrix} >
\]
\[
E v_E ( -\frac{1}{2} + \frac{1}{2} \cdot \sqrt{3}i) = ker \begin{pmatrix}
\frac{3}{2} - \frac{1}{2} \cdot \sqrt{3}i &-2&1 \\
 0&\frac{3}{2} - \frac{1}{2} \cdot \sqrt{3}i&3\\
 0&-1&-\frac{3}{2} - \frac{1}{2} \cdot \sqrt{3}i
 \end{pmatrix}  = ker \begin{pmatrix}
\frac{3}{2} - \frac{1}{2} \cdot \sqrt{3}i &-2&1 \\
 0&\frac{3}{2} - \frac{1}{2} \cdot \sqrt{3}i&3\\
 0&-1&0
 \end{pmatrix} 
\]
\begin{align*}
    \Rightarrow 0 & =  x_2  \cdot (\frac{3}{2} - \frac{1}{2} \cdot \sqrt{3}i) - 2x_2 + x_3 & I \\
    0  & = x_2  \cdot (\frac{3}{2} - \frac{1}{2} \cdot \sqrt{3}i) + 3x_3 & II \\
    \Leftrightarow x_3 & = -x_2 (\frac{1}{2} - \frac{1}{6} \sqrt{3}i) 
\end{align*}
Setze II in I ein
\begin{align*}
   0 & =  x_1 (\frac{3}{2}-\frac{1}{2}\sqrt{3}i) -2x_2 - x_2 \cdot (\frac{1}{2} - \frac{1}{6}\sqrt{3}i) \\
   x_1 & = x_2 \frac{(\frac{5}{2}- \frac{1}{6} \sqrt{3}i)}{(\frac{3}{2}-\frac{1}{2}\sqrt{3}i)} 
\end{align*}
\[
E v_E (\lambda_2) = <
\begin{pmatrix}
x_2 \frac{(\frac{5}{2}- \frac{1}{6} \sqrt{3}i)}{(\frac{3}{2}-\frac{1}{2}\sqrt{3}i)} \\ 1 \\ -\frac{1}{2} + \frac{1}{6}\sqrt{3}i
\end{pmatrix}>
\]
\[
E v_E ( -\frac{1}{2} - \frac{1}{2} \cdot \sqrt{3}i) = ker
\]
Aus Aufgabe 3.2 folgt 
\[
E v_E (\lambda_3) = \overline{E v_E(\lambda_2)}
\]
Frage mit den Auffassung: \\
Reeller Eigenraum: $ < \begin{pmatrix} 1 \\ 0 \\ 0 \end{pmatrix} >$ \\
Reeller Eigenwert: $\lambda_1 = 1$
\section*{Besprechung 3.3}
Bilder (Induktion und Laplace auf die erste Zeile anwenden) 
\section{Aufgabe 5.1}
Das wäre bereits Jordan Normalform: 
\[
\begin{pmatrix}
5 & 1 & 0 \\
0 & 5 & 0 \\
0 & 0 & 6
\end{pmatrix}
\]
JNF existiert nur, wenn das charakteristische Polynom in Linearfaktoren zerfällt. A hat JN  J wenn exists S $\in$ Sln(K)???  \\
$S^{-1} A S = J $
\subsection*{(a)}
Man zeigt sie es nilpotent, wenn man zeigt es gibt ein k mit dem das klappt 
\[
A = \begin{pmatrix}
1& -1& 1\\0 & 0 & 2 \\ 1&-1&-1
\end{pmatrix}
\]
\[
A^2 = \begin{pmatrix}
1& -1& 1\\0 & 0 & 2 \\ 1&-1&-1
\end{pmatrix} \cdot \begin{pmatrix}
1& -1& 1\\0 & 0 & 2 \\ 1&-1&-1
\end{pmatrix} = \begin{pmatrix}
2 & -2 & -2 \\ 2 & -2 & -2 \\ 0 & 0 & 0 
\end{pmatrix}
\]
\[
A^3 = \begin{pmatrix}
2 & -2 & -2 \\ 2 & -2 & -2 \\ 0 & 0 & 0 
\end{pmatrix} \cdot \begin{pmatrix}
1& -1& 1\\0 & 0 & 2 \\ 1&-1&-1
\end{pmatrix} = 0
\] 
JordanBasis: Man will Vektore v1, v2, v3 finden bei denen v1 auf v2 geschickt wird und v2 auf v3 und v3 auf 0 und wenn die eine Basis bilden ist irgendwas. \\
Solche Ketten von Elementen möchte man finden. Es ginge sowohl $v1 \mapsto v2 \mapsto v3 \mapsto$ als auch $v1 \mapsto v2 \mapsto 0, \quad v3 \mapsto 0$
\subsection*{(b)}
Was ist ein Element, das nicht im Kern von $A^2$ ist? setze \\
v3 = $\begin{pmatrix}
1\\0\\0
\end{pmatrix}$
\[v1 = A^2 \cdot v3 = \begin{pmatrix}
2\\2\\0
\end{pmatrix}\]
\[v2 = A \cdot v3 = \begin{pmatrix}
1\\0\\1
\end{pmatrix}\]
\subsection*{(c)}
S ist die geeignete Basiswechselmatrix
\[ S = (v1\ v2\ v3) = \begin{pmatrix} 
2&1&1\\ 2&0&0\\ 0&1&0
\end{pmatrix}
\Rightarrow S^{-1}AS = \begin{pmatrix}
0&1&0\\0&0&1\\0&0&0 
\end{pmatrix}
\]

\section{Aufgabe 5.2}
\subsection*{(a)}
\subsection*{(b)}



\end{document}